\input{../tex-inputs/latex-leseansicht-vorspann}

\section[Arthur Schnitzler an Rainer Maria Rilke, 4. 7. 1901]{L04343 Arthur Schnitzler an Rainer Maria Rilke, 4. 7. 1901}
\nopagebreak\mylabel{L04343v}
\rehead{ }\normalsize\beginnumbering\briefempfaengerindex{Rilke, Rainer Maria@\textsc{Rilke, Rainer Maria}!zzzSchnitzler, Arthur@\emph{von Arthur Schnitzler}!1901-07-044@{4. 7. 1901}|(be}
\toendnotes[C]{\smallbreak\pagebreak[2]}
\correspDesc{Versand  durch Arthur Schnitzler am 4. 7. 1901 in St. Anton am Arlberg
\newline{}Erhalt  durch Rainer Maria Rilke im Zeitraum [5. 7. 1901
                  – 9. 7. 1901?] in Wien}\toendnotes[C]{\smallbreak}
\Standort{DLA, A:Schnitzler, HS.2022.81.}
\physDesc{Brief, 1 Blatt, 3 Seiten, 1020 Zeichen
\newline{}Handschrift: schwarze Tinte, deutsche Kurrent}
\buchAbdrucke{\weitereDrucke{1) \emph{Rainer Maria Rilke und Arthur Schnitzler. Ihr Briefwechsel.}. Mit Anmerkungen herausgegeben von Heinrich Schnitzler In: \emph{Wort und Wahrheit}, Jg. 13, Nr. 4, April 1958, S. 282–298, hier S. 287–288.} \weitereDrucke{2) Arthur Schnitzler: \emph{Briefe 1875–1912}. Herausgegeben von Therese Nickl und Heinrich Schnitzler. Frankfurt am Main: \emph{S. Fischer} 1981, S. 438–439.} }\toendnotes[C]{\smallbreak}
\pstart{}{\pb}Verehrter Herr Rielke,\pend\vspace{0.5em}
\pstart
           Sie beſchämen mich – ich habe Ihnen \introOben{}nur\introOben{} zwei Bücher\pwindex{Schnitzler, Arthur 15. 5. 1862 Wien – 21. 10. 1931 ebd.@\textsc{Schnitzler, Arthur} (15. 5. 1862 Wien – 21. 10. 1931 ebd.)!Schleier der Beatrice. Schauspiel in fünf Akten@\strich\emph{Der Schleier der Beatrice. Schauspiel in fünf Akten}|pwv}\pwindex{Schnitzler, Arthur 15. 5. 1862 Wien – 21. 10. 1931 ebd.@\textsc{Schnitzler, Arthur} (15. 5. 1862 Wien – 21. 10. 1931 ebd.)!Lieutenant Gustl. Novelle@\strich\emph{Lieutenant Gustl. Novelle}|pwv} geſchickt – und erhalte
               von Ihnen zwei \label{K_L04343-1v}\edtext{Briefe}{\lemma{\textnormal{\emph{Briefe}}}\Cendnote{\textnormal{{XXXX ref} und {XXXX ref}.}}}\label{K_L04343-1}, von denen einer herzlicher und
               liebenswürd\textcolor{gray}{ger} iſt als der andre. Ihre \label{K_L04343-11v}\edtext{Flugblätter\pwindex{Frühling. Moderne Flugblätter. 4. Heft@\emph{Frühling. Moderne Flugblätter. 4. Heft}|pwv}\pwindex{Rilke, Rainer Maria 4.\,12.\,1875 Prag – 29.\,12.\,1926 Montreux@\textsc{Rilke, Rainer Maria} (4.\,12.\,1875 Prag – 29.\,12.\,1926 Montreux)!Avalun. [3. Heft: Rainer Maria Rilke]@\strich\emph{Avalun. [3. Heft: Rainer Maria Rilke]}|pwv}}{\lemma{\textnormal{\emph{Flugblätter}}}\Cendnote{\textnormal{Vgl. XXXX Auszeichnungsfehler: Dokument L04344 nicht gefunden.}}}\label{K_L04343-11} hab ich noch nicht; ich find{ }ſie wohl in Wien\oindex{Wien@\textbf{Wien}|pw}, und freue mich im voraus. Augenblicklich befind ich mich
               auf Reiſen, war in Salzburg\oindex{Salzburg@\textbf{Salzburg}|pw}, Innsbruck\oindex{Innsbruck@\textbf{Innsbruck}|pw}, wo es von hier aus hin geht, weiſs {\pb}ich noch nicht zu{ }ſagen.\pend
           
\pstart
           Ihre Worte über die \textsc{Beatrice\pwindex{Schnitzler, Arthur 15. 5. 1862 Wien – 21. 10. 1931 ebd.@\textsc{Schnitzler, Arthur} (15. 5. 1862 Wien – 21. 10. 1931 ebd.)!Schleier der Beatrice. Schauspiel in fünf Akten@\strich\emph{Der Schleier der Beatrice. Schauspiel in fünf Akten}|pw}} haben mich ganz beſonders gefreut. Niemals hab ich mit{ }ſoſehr danach geſehnt,
                  e\textcolor{gray}{t}was von mir wirklich gut geſpielt zu{ }ſehn. Leider hab ich bis
               jetzt nur die{ }ſchwächliche Breslauer\oindex{Breslau@\textbf{Breslau}|pw}{ }Aufführung\pwindex{Schnitzler, Arthur 15. 5. 1862 Wien – 21. 10. 1931 ebd.@\textsc{Schnitzler, Arthur} (15. 5. 1862 Wien – 21. 10. 1931 ebd.)!Schleier der Beatrice. Schauspiel in fünf Akten@\strich\emph{Der Schleier der Beatrice. Schauspiel in fünf Akten}|pwv}\eventindex{Lobe-Theater@\textbf{Lobe-Theater}!Uraufführung von Der Schleier der Beatrice, 1.12.1900@Uraufführung von Der Schleier der Beatrice, 1.12.1900|pwv} erlebt. Wie{ }ſich das einzige Theater\orgindex{Burgtheater@Burgtheater|pwv}, das das Stück\pwindex{Schnitzler, Arthur 15. 5. 1862 Wien – 21. 10. 1931 ebd.@\textsc{Schnitzler, Arthur} (15. 5. 1862 Wien – 21. 10. 1931 ebd.)!Schleier der Beatrice. Schauspiel in fünf Akten@\strich\emph{Der Schleier der Beatrice. Schauspiel in fünf Akten}|pwv}{ }ſpielen konnte, b\textcolor{gray}{e}tragen hat, iſt Ihnen ja
               wahrſcheinlich bekannt. Aber Sie wiſſen ja: ich bin in das Kaſtl mit der Aufſchrift
                  »Liebelei\pwindex{Schnitzler, Arthur 15. 5. 1862 Wien – 21. 10. 1931 ebd.@\textsc{Schnitzler, Arthur} (15. 5. 1862 Wien – 21. 10. 1931 ebd.)!Liebelei. Schauspiel in drei Akten@\strich\emph{Liebelei. Schauspiel in drei Akten}|pw}« hineingerathen; die Kritiker haben
               das {\pb}nicht gern, we{\geminationn} die
               Taferln gewechſelt werden.\pend
           
\pstart
           Ich hoffe es geht Ihnen{ }ſo wohl als ichs Ihnen wünſche.\pend
           
\pstart
           Mit herzlichem Gruſs{\\[\baselineskip]}Ihr aufrichtig ergebner{\\[\baselineskip]}\spacefill\mbox{Arthur Schnitzler}\pend
           \leftskip=0em{}
\pstart
           \noindent{}\textsc{St Anton Arlberg}\oindex{St. Anton am Arlberg@\textbf{St. Anton am Arlberg}|pw}\pend
           
\pstart
           4. 7. 1901\pend
           \selectlanguage{ngerman}\endnumbering\briefempfaengerindex{Rilke, Rainer Maria@\textsc{Rilke, Rainer Maria}!zzzSchnitzler, Arthur@\emph{von Arthur Schnitzler}!1901-07-044@{4. 7. 1901}|)be}\mylabel{L04343h}
\begin{anhang}
\end{anhang}\newcommand{\dateiname}{L04343}\newcommand{\titel}{Arthur Schnitzler an Rainer Maria Rilke, 4. 7. 1901}\newcommand{\editorInnen}{Herausgegeben von Jahnke, SelmaMüller, Martin Anton}\input{../tex-inputs/latex-leseansicht-abspann}
